\section{Likelihood cancellation and the forced binary subtree}\label{sec:coefficients}

Fix the block of Proposition~\ref{prop:block}.  For every supported exact value $\mathfrak d$ of the report $D_0$ from Section~\ref{sec:metadata}, write $\E_{\mathfrak d}$ for expectation over the remaining hidden bits with this report fixed.  This expectation includes failed layouts; it is not conditioned on $\mathcal F$.  The report space is countable, so its supported values carry its full law.

Let $\mathcal L$ be the leaf set and use the conditional reference
\[
 Q=\mu_{i-1}^{\otimes\mathcal L}.
\]
For $v\in\mathcal L$, let $Z_v$ be the old likelihood score at that coordinate.  Under $Q$ these variables are independent and centered, with
$\|Z_v\|_{r_i}^2\le N^{-a_0d}$ by \eqref{eq:old-score}.  Conditional on the bits, the leaf of type $-1$ has density $1+Z_v$ and the leaf of type $+1$ has density one.  Thus the two conditional surrogate densities are
\begin{equation}\label{eq:conditional-densities}
 G_\pm=\E_{\mathfrak d}\prod_{v\in\mathcal L}
       \left(1+\frac{1-\tau_v^\pm}{2}Z_v\right).
\end{equation}
For every supported $\mathfrak d$, these densities integrate to one because they are mixtures of product probability laws.  The superscript on $\tau$ records the root type.

\begin{lemma}[Conditional score gain]\label{lem:tree-score}
For every exact typical metadata value $\mathfrak d$, the densities in \eqref{eq:conditional-densities} satisfy
\[
 \|G_+-G_-\|_{L^{r_i}(Q)}^2\le N^{-d'-.01}.
\]
The bound is uniform in $\mathfrak d$ and the block level $i$.
\end{lemma}

We now fix a typical $\mathfrak d$ for the proof of the lemma.

\subsection{High degrees from the input score}

For fixed bits, every level-one window has at most $L$ negative leaves, because its negative types can occur only on the support of its motif.  There are $\prod_{j=2}^bT_j$ such windows.  This remains true on failed layouts, so the number of factors $1+Z_v$ in either fixed-bit density is at most
\[
 L\prod_{j=2}^bT_j.
\]
The margin between the old score and this number of possible positions is
\begin{equation}\label{eq:score-gap}
 \lambda_d:=a_0d-\sum_{j=2}^b\kappa_j>8g_B(d+1),
 \qquad g_B:=2^{-4B-10}.
\end{equation}
Here is the promised use of the two depth phases.  Since
$\sum_{j=2}^b\kappa_j<1-2^{1-b}$ and $a_0\ge2^{-b}-\epsilon$, for $b=2$ and $d\ge2.20$,
\[
 \lambda_d>d/4-1/2-\epsilon d.
\]
For $b=B$ after the preliminary phase, $d\ge2^B+2.20$ and
\[
 \lambda_d>2^{-B}(d-2^B+2)-\epsilon d.
\]
In each case the positive part is a fixed positive fraction, depending only on $B$, of $d+1$.  The choices of $g_B$ and $\epsilon$ leave the bound in \eqref{eq:score-gap} uniformly.

At fixed bits, Lemma~\ref{lem:moment} bounds the squared $L^{r_i}(Q)$ norm of the homogeneous degree-$r$ part of either density by
\[
 (C_1r_i)^{2r}N^{-a_0dr}
       \binom{L\prod_{j=2}^bT_j}{r}
 \le \bigl(N^{-\lambda_d+o(\epsilon)}\bigr)^r.
\]
Convexity permits averaging over the bits at this bound.  Geometric summation of the norms, using \eqref{eq:score-gap}, gives
\begin{equation}\label{eq:tree_high}
 \|G_\pm-1\|_{r_i}\le N^{-g_B},\qquad
 \|\text{degree }\ge r_*\text{ part of }G_\pm\|_{r_i}^2
       \le N^{-d'-1}.
\end{equation}
For the second inequality, $r_*g_B(d+1)>d'+3$ leaves the required margin.  Only the fixed degrees $1\le r<r_*$ remain.

\subsection{A fixed coefficient and its parity terms}

Fix such a degree $r$ and a leaf subset $\mathcal S$ of size $r$.  Its coefficient in $G_+-G_-$ is
\begin{equation}\label{eq:fixed-coefficient}
 c(\mathcal S)=\E_{\mathfrak d}\left[
       \prod_{v\in\mathcal S}\frac{1-\tau_v^+}{2}
       -\prod_{v\in\mathcal S}\frac{1-\tau_v^-}{2}\right].
\end{equation}
Lemma~\ref{lem:moment} reduces the squared norm of this degree to
\begin{equation}\label{eq:degree-reduction}
 \left\|\sum_{|\mathcal S|=r}c(\mathcal S)
       \prod_{v\in\mathcal S}Z_v\right\|_{r_i}^2
 \le (C_1r_i)^{2r}N^{-a_0dr}
       \sum_{|\mathcal S|=r}|c(\mathcal S)|^2.
\end{equation}
To estimate the right side, expand the integrand for fixed bits and fixed relative leaf labels, and only then average each chosen term.  Call a vertex occupied if it is an ancestor of a leaf in $\mathcal S$, including that leaf itself.

For a vertex $I$ and one of its child labels $h$, write
\[
 \xi_{I,h}=\begin{cases}h-X,&I=\varnothing,\\
 h-w_j+Y_I,&I\text{ at level }j<b,
 \end{cases}
 \qquad \chi_{Ih}=\chi(\xi_{I,h}),\quad
 \sigma_{Ih}=\sigma_{\xi_{I,h}}.
\]
For an occupied vertex $I$, let $F_I(t)$ be the contribution of the selected leaves below it when its incoming type is $t\in\{-1,1\}$.  At a level-one vertex with $m_I$ selected children, the type rule gives the already-cancelled expression
\begin{equation}\label{eq:coefficient-base}
 F_I(t)=2^{-m_I}\prod_{h:\,Ih\in\mathcal S}
       (\chi_{Ih}-\sigma_{Ih}t).
\end{equation}
In particular every selected leaf contributes a motif factor.  At a higher occupied vertex, the selected-leaf product splits across child subtrees, giving the algebraic recursion
\begin{equation}\label{eq:coefficient-recursion}
 F_I(t)=\prod_{h:\,Ih\text{ occupied}}
       F_{Ih}(1-\chi_{Ih}+\sigma_{Ih}t),\qquad t\in\{-1,1\}.
\end{equation}
The substituted type is again $+1$ or $-1$.  No probabilistic independence of the child functions is asserted here; this is the factorization of one product of leaf indicators at fixed bits.

Every function on $\{-1,1\}$ has an even and an odd part.  To retain the individual motif choices, expand \eqref{eq:coefficient-base}--\eqref{eq:coefficient-recursion} into terms $M_I t^{p_I}$, reducing powers by $t^2=1$, where $p_I\in\{0,1\}$.  The following rules describe each chosen term.
\begin{itemize}
\item At level one choose $\chi_{Ih}$ or $-\sigma_{Ih}t$ from every selected leaf factor.  Every such leaf is called active, and $p_I$ is the parity of the number of signed factors.
\item At a higher vertex choose one term $M_{Ih}\tau_{Ih}^{p_{Ih}}$ from each occupied child.  If $p_{Ih}=0$, its contribution is $M_{Ih}$.  If $p_{Ih}=1$, substitute the type rule and choose one of $1$, $-\chi_{Ih}$, or $\sigma_{Ih}t$ in
$M_{Ih}(1-\chi_{Ih}+\sigma_{Ih}t)$.  The child is active at $I$ when the second or third summand is chosen.  Set $p_I$ to the parity of the signed summands chosen at $I$.
\end{itemize}
Odd and even are roles of these chosen terms, not intrinsic properties of the vertex or of its full coefficient function.  At the root, the difference $F_\varnothing(+1)-F_\varnothing(-1)$ keeps exactly the odd terms and gives each a factor two.  Apart from bounded scalar factors, a chosen term in \eqref{eq:fixed-coefficient} is therefore the expectation of the product of its active $\sigma$ and $\chi$ factors.

For fixed $r$ and $b$, the number of ordered occupied-tree shapes and expansion choices is bounded solely in terms of $r,b$.  A shape specifies adjacency and sibling order, not the numerical child labels.  We may fix a shape and its choices, bound the squared coefficient sum over the labels, and combine these finitely many bounds by Cauchy--Schwarz.

We first split that label sum into central and noncentral assignments.  If an active child of a level-$j$ vertex has label outside $[T_j/3,2T_j/3]$, a nonzero factor requires $|Y_I|>T_j/8$ for all large instances, since $w_j=T_j/2$ and $L=o(T_j)$; at the root it is impossible by \eqref{eq:outer-anchor}.  Hence such a term has conditional multiplier at most
$\Prob(\mathcal F^c\mid D_0=\mathfrak d)\le\exp(-N^{\epsilon/4})$ in absolute value.  There are at most $O_B(N^r)$ label assignments by \eqref{eq:power_gap}, so their contribution to the right side of \eqref{eq:degree-reduction} is at most $N^{-d'-1}$ eventually.  From now on all active child labels in the sum lie in their central intervals.  This is a restriction of deterministic label assignments, not conditioning of the bit law.  In particular, two forced children of a level-$(j+1)$ parent will lie in
\[
 [T_{j+1}/3,2T_{j+1}/3]\subset(T_{j+1}/4,3T_{j+1}/4],
\]
the interval where their offsets share the global sum in \eqref{eq:global-plateau}.

\subsection{Terms with a singleton odd role}

Suppose a chosen odd term at a vertex $I$ has exactly one active child.  The active factor must be signed, since the number of signed factors is odd.  All other local choices at that vertex contribute one or no factor.

We need the precise locality of its relative argument.  After the relative IDs and expansion choices are fixed, the local sum $H_I$ occurs among the surrogate arguments $\xi_{J,h}$ only when $J=I$.  Equation~\eqref{eq:prefix-offset} shows that the local pre/post deviations cancel before a later sibling part.  Lower relative arguments use the fresh mechanisms of their own groups.  The physical absolute positions of the descendant subtree can move with $H_I$, but those positions are unreported and do not occur in its lower relative arguments.  Likewise $X$ occurs only in root arguments.  This statement does not require independence of child terms, which may share lower global bits.

Condition on every orientation bit except the hidden local bits at $I$; for the root, keep only its hidden outer bits unconditioned.  Every other active factor is then fixed, while the singleton signed factor is convolved with a translated or reflected fair binomial.  Its number of trials is at least $\omega^2t_j/2$ for a level-$j$ vertex with $j<b$, and at least $\omega^2J/2$ at the root.  Equation~\eqref{eq:sig_conv} bounds the conditional multiplier by
\[
 \exp(O((s_*+1)^2+(s_*+1)\Lambda))
 \begin{cases}
 t_j^{-(s_*+1)/2},&1\le j<b,\\
 J^{-(s_*+1)/2},&j=b.
 \end{cases}
\]
The estimate remains valid after averaging the other bits.  Typicality is already a property of the fixed report, and no fit conditioning has been used.

The lower bounds for $J,t_j$ give
\[
 J^{-(s_*+1)},\ t_j^{-(s_*+1)}
 \le e^{O_B((s_*+1)(D+1)\Lambda)}
       N^{-\epsilon(s_*+1)}.
\]
Since $\epsilon s_*\ge100(d_D+r_*+B+10)$, this absorbs the $O_B(N^r)$ label assignments and the squared moment factor $(C_1r_i)^{2r}N^{-a_0dr}$.  The resulting contribution to the squared degree-$r$ norm is at most $N^{-d'-1}$ eventually.

\subsection{The forced subtree and extra labels}

In a remaining odd root term, every odd nonleaf has at least two active children.  At a higher vertex those children carry odd chosen terms; at level one they are selected leaves.  Starting at the root, choose the first two active children in sibling order and repeat this choice at each selected odd child.  This forces a complete binary subtree with $2^b$ leaves.  The choice depends on the shape and expansion roles, not on the numerical labels.  Thus $r\ge2^b$; put $e=r-2^b$.

For each forced nonleaf $I$, let $h_I<h'_I$ be its two forced child labels, calling the first the anchor.  A nonzero term requires
$0<h'_I-h_I<L$, because both active arguments belong to the same translate of $\mathcal I$.  Define the nonnegative anchor kernel
\[
 f=\prod_{I\text{ forced nonleaf}}\chi(\xi_{I,h_I}).
\]
Every extra, nonforced selected leaf still contributes an absolute $\chi$ constraint by \eqref{eq:coefficient-base}.  After discarding the other constraints, the absolute multiplier is at most a fixed scalar times
\begin{equation}\label{eq:anchor-extra-bound}
 \E_{\mathfrak d}\left[f\prod_{\text{extra leaves }Ih}
       \chi(\xi_{I,h})\right].
\end{equation}
The expectation has no additional fit indicator.  The central and close-label restrictions concern the fixed labels being summed.

Figure~\ref{fig:forced-subtree-jitter} illustrates the forced selection and the shared global offset produced by \eqref{eq:prefix-offset}--\eqref{eq:global-plateau}.
\begin{figure}[p]
\centering
\begin{tikzpicture}[x=1cm,y=1cm,font=\small,>=Latex,
  odd/.style={circle,draw=linkblue,fill=linkblue!10,inner sep=2pt,
              minimum size=0.62cm,font=\scriptsize},
  leaf/.style={circle,fill=linkblue,inner sep=0pt,minimum size=4.5pt},
  box/.style={draw=linkblue,rounded corners=2pt,fill=linkblue!5,
              align=center,inner sep=4pt}]
  \node[anchor=west,font=\small\bfseries] at (0,0.35)
    {(a) Chosen terms: singleton smoothing and forced selection};

  \node[odd] (single-root) at (1.8,-0.25) {odd};
  \node[leaf] (single-child) at (1.8,-1.10) {};
  \draw[linkblue,thick] (single-root)--(single-child)
    node[midway,right=3pt] {$\sigma$};
  \node[align=center] at (1.8,-1.50) {one active child};
  \node[box,text width=3.8cm] (single-result) at (1.8,-3.00)
    {signed smoothing\\[1pt]
     negligible term};
  \draw[->,linkblue] (1.8,-1.75)--(single-result.north)
    node[midway,right=3pt,font=\scriptsize] {average hidden bits};

  \node[odd] (root) at (9.6,-0.25) {odd};
  \node[anchor=west] at (10.07,-0.25) {root};
  \node[odd] (a) at (7.6,-1.10) {odd};
  \node[odd] (b) at (11.6,-1.10) {odd};
  \node[odd] (aa) at (6.6,-1.95) {odd};
  \node[odd] (ab) at (8.6,-1.95) {odd};
  \node[odd] (ba) at (10.6,-1.95) {odd};
  \node[odd] (bb) at (12.6,-1.95) {odd};
  \foreach \parent/\child in {root/a,root/b,a/aa,a/ab,b/ba,b/bb}
    \draw[linkblue,thick] (\parent)--(\child);
  \foreach \parent/\xx/\yy in
    {aa/6.1/7.1,ab/8.1/9.1,ba/10.1/11.1,bb/12.1/13.1} {
    \node[leaf] (left-\parent) at (\xx,-2.80) {};
    \node[leaf] (right-\parent) at (\yy,-2.80) {};
    \draw[linkblue,thick] (\parent)--(left-\parent);
    \draw[linkblue,thick] (\parent)--(right-\parent);
  }
  \draw[decorate,decoration={brace,mirror,amplitude=4pt},linkblue]
    (5.98,-3.00)--(13.22,-3.00)
    node[midway,below=5pt,text=black] {$2^b$ selected leaves ($b=3$ shown)};

  \draw[black!20] (0,-4.00)--(14.2,-4.00);
  \node[anchor=west,font=\small\bfseries] at (0,-4.35)
    {(b) Two forced siblings share a global offset};
  \node[anchor=west] at (0,-5.10)
    {Fixed $G$ at level $j+1$; distinct central siblings $I_1=Gu_1$, $I_2=Gu_2$};
  \fill[black!4] (0.8,-6.40) rectangle (13.6,-6.85);
  \fill[linkblue!15] (4.0,-6.40) rectangle (10.4,-6.85);
  \draw[black!50] (0.8,-6.40) rectangle (13.6,-6.85);
  \draw[linkblue,dashed] (4.0,-6.30)--(4.0,-6.95);
  \draw[linkblue,dashed] (10.4,-6.30)--(10.4,-6.95);
  \node[anchor=north west] at (0.8,-6.90) {$1$};
  \node[anchor=north east] at (13.6,-6.90) {$T_{j+1}$};
  \node[anchor=south] at (4.0,-6.15) {$T_{j+1}/4$};
  \node[anchor=south] at (10.4,-6.15) {$3T_{j+1}/4$};
  \node[leaf] at (6.05,-6.625) {};
  \node[leaf] at (8.35,-6.625) {};
  \node[anchor=south] at (6.05,-6.32) {$I_1$};
  \node[anchor=south] at (8.35,-6.32) {$I_2$};
  \draw[decorate,decoration={brace,mirror,amplitude=4pt},linkblue]
    (4.02,-7.00)--(10.38,-7.00)
    node[midway,below=5pt,text=black] {middle-half plateau: $V_I=V_G$};

  \node[box,minimum width=2.1cm] (global) at (7.2,-8.35) {shared $V_G$};
  \node[box,fill=white,minimum width=2.1cm] (local1) at (1.8,-8.35)
    {local $H_{I_1}$};
  \node[box,fill=white,minimum width=2.1cm] (local2) at (12.6,-8.35)
    {local $H_{I_2}$};
  \node[box,minimum width=4.1cm] (offset1) at (3.6,-9.65)
    {$Y_{I_1}=V_G+H_{I_1}$};
  \node[box,minimum width=4.1cm] (offset2) at (10.8,-9.65)
    {$Y_{I_2}=V_G+H_{I_2}$};
  \draw[->,linkblue,thick] (global.south)--([xshift=8mm]offset1.north);
  \draw[->,linkblue,thick] (global.south)--([xshift=-8mm]offset2.north);
  \draw[->,black!65] (local1.south)--([xshift=-8mm]offset1.north);
  \draw[->,black!65] (local2.south)--([xshift=8mm]offset2.north);
  \node[align=center] at (7.2,-10.50)
    {$H_{I_1},H_{I_2},V_G$ are independent given $D_0=\mathfrak d$.};
\end{tikzpicture}
\caption{The labels ``odd'' in panel (a) describe chosen-term roles.
A singleton odd term, including at the root, is negligible after signed
smoothing.  For a remaining odd root term, selecting the first two active
children recursively forces a binary subtree with $2^b$ selected leaves
and hence $r\ge2^b$ ($b=3$ shown).  The selection is fixed by the term roles
and sibling order.  Panel (b) illustrates
\eqref{eq:prefix-offset}--\eqref{eq:global-plateau} for fixed distinct
central sibling IDs $u_1<u_2$ in one group $G$.  Conditional on the exact
typical report $D_0=\mathfrak d$, the local sums $H_{I_1},H_{I_2}$ and the
global sum $V_G$ are independent; the total offsets $Y_{I_1},Y_{I_2}$ share
$V_G$ and are estimated jointly.}
\label{fig:forced-subtree-jitter}
\end{figure}

\clearpage

First sum the extra labels so that only the forced anchor expectation remains.  Fix the forced labels and all nonleaf labels.  For an assignment $\mathbf h$ of the $e$ extra leaf labels, let $a_{\mathbf h}$ be the nonnegative expectation in \eqref{eq:anchor-extra-bound}.  At fixed bits each extra leaf has at most $L$ labels satisfying its constraint, even if several constraints share bits.  Hence
\[
 \sum_{\mathbf h}a_{\mathbf h}\le L^e\E_{\mathfrak d}f,
 \qquad
 \sum_{\mathbf h}a_{\mathbf h}^2
 \le\left(\sum_{\mathbf h}a_{\mathbf h}\right)^2
 \le L^{2e}(\E_{\mathfrak d}f)^2.
\]
This avoids a free factor $T_1$ per extra leaf.  At each positive nonroot level, there are at most $e$ occupied nonforced vertices: their descendant sets are disjoint at that level, and each contains an extra leaf and no forced leaf.  Summing their child labels costs at most
$(\prod_{j=2}^bT_j)^e$.  With the forced labels fixed, $f$ does not depend on which nonforced labels were selected, because the complete tree and every group mechanism were fixed before choosing $\mathcal S$.

\subsection{The fixed-label anchor factorization}

It remains to sum $(\E_{\mathfrak d}f)^2$ over the forced labels.  Call a forced sibling pair valid when its two labels are central and their positive difference is less than $L$.  Dropping any further order restrictions involving nonforced children can only increase this nonnegative sum.  Define the outer anchor convolution
\[
 K_{\rm out}^{\mathfrak d}(z)=\E_{\mathfrak d}\chi(z-X).
\]
Its $\ell^1$ mass is $L$, and its supremum is at most $O(L\omega^{-1}J^{-1/2})$ by the hidden root count.  Including at most $2L$ choices for the close root child, its squared sum costs
\begin{equation}\label{eq:outer-anchor-cost}
 \sum_{\substack{u_1,u_2\text{ valid root}\\\text{forced labels}}}
       K_{\rm out}^{\mathfrak d}(u_1)^2
 \le \Gamma_{\rm out},\qquad
 \Gamma_{\rm out}=O(L^3\omega^{-1}J^{-1/2}).
\end{equation}

For a fixed forced labeling, let $\mathcal G_j$ be its forced level-$(j+1)$ parents.  Each $G\in\mathcal G_j$ has two forced level-$j$ children $Gu_1(G),Gu_2(G)$ with fixed distinct labels $u_1(G)<u_2(G)$ in the middle-half interval.  For any such fixed $G$ and sibling labels $u_1<u_2$, define their joint anchor kernel by
\begin{equation}\label{eq:joint-anchor}
 \mathcal H_{j,G;u_1,u_2}^{\mathfrak d}(z_1,z_2)
 =\E_{\mathfrak d}\left[
   \chi(z_1-w_j+Y_{Gu_1})\chi(z_2-w_j+Y_{Gu_2})\right].
\end{equation}
The expectation here uses that group's global bits and the two local collections under the fixed report.  The actual $G,u_1,u_2$ select those collections; only $z_1,z_2$ are variables in the norm below.  The factor can vary with all of these fixed data.

Independence of the distinct group mechanisms conditional on $D_0=\mathfrak d$ now gives, for this fixed forced labeling,
\begin{equation}\label{eq:anchor-factorization}
 \E_{\mathfrak d}f
 =K_{\rm out}^{\mathfrak d}(h_\varnothing)
   \prod_{j=1}^{b-1}\prod_{G\in\mathcal G_j}
   \mathcal H_{j,G;u_1(G),u_2(G)}^{\mathfrak d}
       (h_{Gu_1(G)},h_{Gu_2(G)}).
\end{equation}
The group sets and sibling IDs in this identity are determined by the fixed upper labels.  In particular $\mathcal G_{b-1}$ contains the root; its interior global/local mechanism is independent of the outer mechanism defining $X$.

For fixed distinct plateau siblings, \eqref{eq:global-plateau} gives
$Y_{Gu_1}=V_G+H_{Gu_1}$ and $Y_{Gu_2}=V_G+H_{Gu_2}$.  Conditional on $\mathfrak d$, the two local hidden sums are independent of each other and of the shared global hidden sum, though their visible shifts and hidden counts can differ.  We have uniformly in the fixed data
\begin{equation}\label{eq:joint-anchor-norms}
\begin{aligned}
 \|\mathcal H_{j,G;u_1,u_2}^{\mathfrak d}\|_{\ell^1(\mathbb Z^2)}&\le L^2,\\
 \|\mathcal H_{j,G;u_1,u_2}^{\mathfrak d}\|_{\ell^\infty(\mathbb Z^2)}
 &\le C_{10}L^2\omega^{-2}(T_{j+1}t_j)^{-1/2}t_j^{-1/2}.
\end{aligned}
\end{equation}
Here $C_{10}$ is absolute.  The first bound follows by summing the two indicators before averaging.  For the supremum, first fix all global bits and bound the local average at $Gu_2$ by $O(L\omega^{-1}t_j^{-1/2})$.  This bound is uniform in the global offset, so discard that factor's constraint.  Average the other factor over its local bits and the global bits.  The latter contain at least a constant times $\omega^2T_{j+1}t_j$ hidden fair trials, giving $O(L\omega^{-1}(T_{j+1}t_j)^{-1/2})$.  Multiplication proves the displayed joint bound without treating the two offsets as independent.

Multiplying the $\ell^1$ and supremum bounds controls the squared $\ell^2$ sum.  Allowing a close child for each anchor costs at most $(2L)^2$.  Thus for fixed valid $G,u_1,u_2$,
\begin{equation}\label{eq:group-anchor-cost}
\begin{aligned}
 \sum_{\substack{z_1,z'_1,z_2,z'_2\in\mathbb Z\\
                  0<z'_1-z_1<L,\ 0<z'_2-z_2<L}}
       \mathcal H_{j,G;u_1,u_2}^{\mathfrak d}(z_1,z_2)^2
 &\le\Gamma_j,\\
 \Gamma_j&=C_{11}L^6\omega^{-2}T_{j+1}^{-1/2}t_j^{-1},
\end{aligned}
\end{equation}
with an absolute $C_{11}$ large enough for every group.

We make the order of the nested sums explicit.  For two fixed valid level-$j$ siblings $Gu_1,Gu_2$, let $\mathcal R_j(G;u_1,u_2)$ be the sum of the squares of the products of all their interior $\mathcal H$ factors, including the factor of $G$, over all valid forced descendant labelings.  Set $\mathcal R_0=1$.  Splitting the descendant labels at their first two pairs gives the exact recurrence
\begin{equation}\label{eq:anchor-sum-recursion}
\begin{aligned}
 \mathcal R_j(G;u_1,u_2)
 &=\sum_{\substack{(z_1,z'_1),(z_2,z'_2)\\
                    \text{valid at level }j}}
 \mathcal H_{j,G;u_1,u_2}^{\mathfrak d}(z_1,z_2)^2
 \prod_{v=1}^2\mathcal R_{j-1}(Gu_v;z_v,z'_v)\\
 &\le \Gamma_j\left(\sup_{G',v_1,v_2\text{ valid}}\mathcal R_{j-1}(G';v_1,v_2)\right)^2.
\end{aligned}
\end{equation}
For the inequality, first use the uniform descendant bound on the valid distinct child pairs.  Only afterward enlarge the remaining close-label sums to at most $2L$ choices and apply \eqref{eq:group-anchor-cost}.  No group estimate is invoked at coincident sibling IDs.  Iterating this inequality from the leaves upward, and finally using \eqref{eq:outer-anchor-cost}, gives
\begin{equation}\label{eq:forced-anchor-sum}
 \sum_{\text{valid forced labels}}(\E_{\mathfrak d}f)^2
 \le \Gamma_{\rm out}\prod_{j=1}^{b-1}\Gamma_j^{2^{b-j-1}}.
\end{equation}
This procedure applies a uniform bound to each actual group after its upper IDs are fixed; it does not replace the varying kernels by one stationary kernel.

\subsection{The exponent gain and closure}

Combine the extra-label reduction, \eqref{eq:forced-anchor-sum}, and the squared moment factor $(C_1r_i)^{2r}N^{-a_0dr}$.  The fixed shape and expansion counts, the powers of $L$ and $\omega^{-1}$, and the positive factor $(C_1r_i)^{2r}$ contribute $N^{o(\epsilon)}$ by Lemma~\ref{lem:scale-margins} and \eqref{eq:block-losses}.  The squared degree contribution of a branching term is therefore at most
\[
 N^{o(\epsilon)}J^{-1/2}
 \prod_{j=1}^{b-1}\bigl[T_{j+1}^{-1/2}t_j^{-1}\bigr]^{2^{b-j-1}}
 N^{-2^ba_0d}N^{-\lambda_de}.
\]
The last factor pays for every extra leaf and is at most one by \eqref{eq:score-gap}.  Using the lower bounds for $J,t_j$, the decay exponent of the remaining factors is at least
\begin{align*}
 &\epsilon/2+\sum_{j=1}^{b-1}2^{b-j-1}(\kappa_{j+1}/2+\epsilon)
       +2^ba_0d-o(\epsilon)\\
 &\qquad\ge d+(b-1)/4
       -\bigl[(2^b-1)d+3b\,2^{b-1}\bigr]\epsilon-o(\epsilon)
       >d'+.03.
\end{align*}
For this calculation, \eqref{eq:local-exponents} gives
$\kappa_{j+1}=2^{j-b}-(6\cdot2^j-2-2^{j-b})\epsilon$.  Its leading contribution at depth $j$ is
$2^{b-j-1}2^{j-b}/2=1/4$, which is the gain from the shared location of each forced sibling group.  The bound $d\le2.20+(B/4)D$ and the chosen $\epsilon$ absorb the displayed corrections and leave the strict margin over $d'=d+(b-1)/4-.05$.

For each $r<r_*$, combine the fixed-shape estimates by Cauchy--Schwarz in \eqref{eq:degree-reduction}.  The noncentral and singleton terms are at most $N^{-d'-1}$ in squared norm, and the branching terms have the margin just proved.  Summing the finitely many degree norms and the tail in \eqref{eq:tree_high} yields
\[
 \|G_+-G_-\|_{r_i}^2\le N^{-d'-.01},
\]
which proves Lemma~\ref{lem:tree-score}.

Apply Lemma~\ref{lem:clipping} to the preliminary laws $\widetilde\mu_i,\widetilde\nu_i$, their common report $D_0$, and the conditional densities \eqref{eq:conditional-densities}.  The typical report has probability $1-o(\Delta)$, \eqref{eq:tree_high} supplies $\rho=g_B$, and Lemma~\ref{lem:tree-score} supplies the density difference.  Conditioning on the common event consisting of typical metadata and $G_+\ge1/2$ gives laws $\mu_i,\nu_i$ with the required domination and score.  Each preliminary law changes by $o(\Delta)$.  The hidden-bit law was used only before this conditioning.

Together with \eqref{eq:projection_tree}, the new conditioning and failed-layout loss is at most $\Delta$ for all large instances.  The inherited loss is at most $(\prod_jT_j)E_{i-1}\le(\prod_jP_j)E_{i-1}=(\ell_i/\ell)E_{i-1}$.  Thus one may take
\[
 E_i\le\Delta+(\ell_i/\ell)E_{i-1},\qquad
 \frac{E_i}{\ell_i}\le\frac{\Delta}{\ell_i}
       +\frac{E_{i-1}}{\ell_{i-1}}\le(i+1)\Delta.
\]
The length and overlap conditions were proved in Section~\ref{sec:block}, so this completes Proposition~\ref{prop:block}.  Starting from Proposition~\ref{prop:initial} and iterating the block proves Proposition~\ref{prop:construction}.
